\section{Metodologia}
\label{chap:metodologia}

\subsection{Análise exploratória e preparação dos dados}
O conjunto de dados \textit{Real Estate Valuation} (UCI) foi organizado em matriz de atributos $\mathbf{X}$ e vetor-alvo $\mathbf{y}$ , correspondente ao preço por unidade de área. A inspeção inicial incluiu estatísticas descritivas, verificação de valores ausentes, não foram identificados dados faltantes, análise de histogramas univariados para avaliar assimetria e presença de caudas, além de uma análise bivariada com base na matriz de correlação entre atributos e alvo. Constatou-se baixa associação linear entre a data de transação e o preço, o que motivou testes preliminares com e sem a utilização dessa variável. Optou-se, contudo, por mantê-la nas análises finais a fim de preservar a comparabilidade com o estudo de referência. 

\subsection{Esquemas de normalização e experimentos de controle}
Foram conduzidos quatro experimentos paralelos para avaliar o efeito do pré-processamento das entradas. Salvo indicação em contrário, apenas $\mathbf{X}$ foi transformada e $\mathbf{y}$ permaneceu na escala original. Os esquemas investigados foram: (i) sem normalização; (ii) padronização por z-score (centralização pela média e escala pelo desvio-padrão); (iii) transformação \textit{min--max} (reescala para $[0,1]$); e (iv) normalização robusta por mediana e intervalo interquartil (IQR). Todas as transformações foram ajustadas exclusivamente nos dados de treino de cada \textit{fold} e aplicadas aos conjuntos de validação e teste, de modo a evitar vazamento de informação. 
O PS emprega ativação sigmoide na saída, cujo suporte é $(0,1)$. Para compatibilizar a faixa do alvo $y\in [8,110]$ com essa saída, adotamos \textbf{normalização \textit{min--max} apenas para $\mathbf{y}$ no treinamento/validação do PS}. As predições do PS foram sempre revertidas para a escala original antes de calcular métricas e gerar gráficos. Nos demais modelos, $\mathbf{y}$ permaneceu na escala original.



\subsection{Protocolo de treinamento, validação e teste}
Os dados foram particionados uma única vez em treino/teste (\SI{66}{\percent}/\SI{34}{\percent}), como no estudo de referência. A seleção de hiperparâmetros ocorreu \emph{exclusivamente} no conjunto de treinamento, por Validação Cruzada (CV) com $k=5$. Para cada configuração, o modelo foi ajustado nos \textit{folds} de treino e avaliado nos \textit{folds} de validação, com as normalizações estimadas apenas nos dados de treino de cada partição. Ao término desse processo, a configuração com melhor desempenho médio em validação foi selecionada. Em seguida, o modelo final foi reestimado utilizando todo o conjunto de treinamento e, por fim, avaliado no conjunto de teste retido. Todas as partições e inicializações foram conduzidas com sementes fixas, assegurando a plena reprodutibilidade dos resultados. Além das médias de CV, agregamos predições \emph{out-of-fold} (OOF) para diagnóstico global de resíduos em treino, sem uso do conjunto de teste.

\subsection{Métricas e procedimentos de avaliação}
A aderência dos modelos foi examinada por meio de gráficos de dispersão entre valores observados e preditos (treinamento e teste), acompanhados dos respectivos coeficientes de correlação linear. 
As métricas primárias de desempenho consideradas foram Hit@10, Hit@20, $R^2$ e RMSE. Adicionalmente, reportamos o RMSE normalizado pelo preço médio (RMSE/Preço).

\paragraph{Definição formal do Hit@$k$.} 
Para cada amostra $i$, dizemos que houve acerto relativo a $k$ se
\begin{equation}
\frac{|y_i - \hat{y}_i|}{y_i} \le \frac{k}{100}.
\end{equation}
O indicador Hit@$k$ é a fração de acertos no conjunto considerado:
\begin{equation}
\text{Hit@}k = \frac{1}{n}\sum_{i=1}^{n} \mathbf{1}\!\left( \frac{|y_i - \hat{y}_i|}{y_i} \le \tfrac{k}{100} \right),
\end{equation}
onde $\mathbf{1}(\cdot)$ é a função indicadora.

Os resíduos de treino dos modelos selecionados foram analisados via histograma e teste de Shapiro--Wilk (nível $\alpha=0{,}05$). Sempre que observada discrepância em relação à gaussianidade (caudas pesadas/assimetria), discutimos implicações para inferência e escolha de métricas.

\subsection{Critério de seleção de modelos}
Para escolher um ``melhor'' por família (MQ, PS, MLP-1C, MLP-2C e \emph{benchmarks}), 
adotamos critério lexicográfico: 
\textbf{Hit@10 (maior)} $\to$ \textbf{Hit@20 (maior)} $\to$ $\mathbf{R^2}$ (maior) $\to$ \textbf{RMSE (menor)}. 
O critério é aplicado uniformemente em todas as buscas e relatórios.

\subsection{Espaços de hiperparâmetros e procedimentos de busca}
A seleção de hiperparâmetros foi conduzida por meio de \textit{random search}, com orçamento de 60 testes, aplicado às redes neurais e ao PS. Para o MQ, empregou-se \textit{grid search}. Em todos os casos, a busca foi realizada com validação cruzada $k=5$, com as normalizações estimadas exclusivamente a partir dos dados de treinamento.


% Please add the following required packages to your document preamble:
% \usepackage{multirow}
% \usepackage{graphicx}
\begin{table}[h]
\caption{Espaços de hiperparâmetros e procedimentos de busca utilizados nos modelos.}
\label{tab:hiperparams-all}
\resizebox{\columnwidth}{!}{%
\begin{tabular}{llll}
\hline
\textbf{Modelo} &
  \textbf{Hiperparâmetros} &
  \textbf{Valores considerados} &
  \textbf{Procedimento de busca} \\ \hline
\multirow{6}{*}{\textbf{MLP-1C}} &
  Neurônios na camada oculta &
  4, 8, 16, 32, 64, 128, 256, 512 &
  \multirow{6}{*}{\begin{tabular}[c]{@{}l@{}}Random Search\\ (60 amostras)\end{tabular}} \\
 & Função de ativação     & \texttt{sigmoid}, \texttt{tanh}, \texttt{relu}, \texttt{leaky\_relu} &  \\
 & Épocas                 & 100, 150, 200, 300                                                   &  \\
 & Taxa de aprendizado    & $10^{-3}, 5\cdot10^{-3}, 10^{-2}$                                    &  \\
 & Normalização           & none, zscore, minmax, iqr                                            &  \\
 & Critério de seleção    & desempenho médio em CV (5 folds)                                     &  \\ \hline
\multirow{6}{*}{\textbf{MLP-2C}} &
  Neurônios na 1ª camada &
  4, 8, 16, 32, 64, 128, 256, 512 &
  \multirow{6}{*}{\begin{tabular}[c]{@{}l@{}}Random Search \\ (60 amostras)\end{tabular}} \\
 & Neurônios na 2ª camada & 4, 8, 16, 32, 64, 128, 256, 512                                      &  \\
 & Função de ativação     & \texttt{sigmoid}, \texttt{tanh}, \texttt{relu}, \texttt{leaky\_relu} &  \\
 & Épocas                 & 100, 150, 200, 300                                                   &  \\
 & Taxa de aprendizado    & $10^{-3}, 5\cdot10^{-3}, 10^{-2}$                                    &  \\
 & Normalização           & none, zscore, minmax, iqr                                            &  \\ \hline
\multirow{4}{*}{\textbf{PS (PS)}} &
  Épocas &
  100, 150, 200, 300 &
  \multirow{4}{*}{\begin{tabular}[c]{@{}l@{}}Random Search \\ (60 amostras)\end{tabular}} \\
 & Taxa de aprendizado    & $10^{-3}, 5\cdot10^{-3}, 10^{-2}$                                    &  \\
 & Normalização           & none, zscore, minmax, iqr                                            &  \\
 & Ativação               & sigmoide (saída compatível com min--max do alvo)                     &  \\ \hline
\multirow{3}{*}{\textbf{MQ (Ridge)}} &
  Parâmetro de regularização ($\lambda$) &
  0.0, $10^{-4}$, $10^{-3}$, $10^{-2}$, $5\cdot10^{-2}$, $10^{-1}$, 0.5, 1.0, 2.0, 5.0, 10.0 &
  \multirow{3}{*}{Grid Search} \\
 & Normalização           & none, zscore, minmax, iqr                                            &  \\
 & Critério de seleção    & desempenho médio em CV (5 folds)                                     &  \\ \hline
\end{tabular}%
}
\end{table}

\subsection{Modelos de referência e comparação}

Como \emph{benchmark}, ajustamos Support Vector Regression (SVR), Random Forest e Gradient Boosting (\texttt{scikit-learn}) sob o mesmo protocolo de particionamento, normalização das entradas e validação cruzada do conjunto de treinamento, com avaliação final no teste retido pelas mesmas métricas. 
Houve uma busca manual dos hiperparâmetros para esses \emph{benchmarks}; os valores foram \textbf{pré-fixados} e mantidos constantes em toda a avaliação.

\begin{table}[h]
\centering
\caption{Modelos e hiperparâmetros selecionados.}
\label{tab:modelos_hiperparam}
\resizebox{\textwidth}{!}{%
\begin{tabular}{@{}lll@{}}
\toprule
\textbf{Modelo} & \textbf{Pré-processamento} & \textbf{Hiperparâmetros} \\ \midrule
Random Forest & \texttt{StandardScaler} &
\begin{tabular}[c]{@{}l@{}}
n\_estimators=100; max\_depth=None; max\_features=1.0; \\
min\_samples\_split=2; min\_samples\_leaf=1; bootstrap=True; \\
criterion='squared\_error'; random\_state=42; n\_jobs=-1
\end{tabular} \\ \midrule

Gradient Boosting & \texttt{RobustScaler} &
\begin{tabular}[c]{@{}l@{}}
n\_estimators=100; learning\_rate=0.1; max\_depth=3; subsample=1.0; \\
loss='squared\_error'; criterion='friedman\_mse'; alpha=0.9; \\
random\_state=42; tol=1e-4; validation\_fraction=0.1
\end{tabular} \\ \midrule

SVR (RBF) & \texttt{RobustScaler} &
\begin{tabular}[c]{@{}l@{}}
kernel='rbf'; C=1.0; epsilon=0.1; gamma='scale'; \\
tol=1e-3; max\_iter=-1; shrinking=True
\end{tabular} \\ \bottomrule
\end{tabular}%
}
\end{table}


\subsection{Implementação e Reprodutibilidade}
\label{sec:implementacao_reprodutibilidade}

Os modelos principais (MQ, PS, MLP-1C e MLP-2C) foram implementados em Python com rotinas matriciais de baixo nível, o que permitiu controle direto sobre inicialização, função de custo e retropropagação. Os \emph{benchmarks} comparativos foram ajustados com \texttt{scikit-learn}. Todo o material necessário para reprodução dos experimentos, incluindo análises exploratórias (EDA), variações de normalização e geração de figuras, encontra-se disponível em repositório público no GitHub:
\begin{center}
    \url{https://github.com/LucasJLBraz/Trabalhos_IC_Aplicada}
\end{center}
As versões de Python e bibliotecas, assim como as sementes aleatórias utilizadas, estão registradas nos notebooks e no arquivo de configuração do projeto.

A Tabela~\ref{tab:estrutura_arquivos_tc1} resume os arquivos e diretórios mais relevantes para a reprodutibilidade do Trabalho~1.

\begin{table}[h!]
\centering
\caption{Arquivos e diretórios relevantes para a reprodutibilidade do Trabalho 1.}
\label{tab:estrutura_arquivos_tc1}
\begin{tabular}{p{5.5cm} p{8cm}}
\toprule
\textbf{Caminho} & \textbf{Descrição} \\
\midrule
\texttt{README.md} & Instruções detalhadas para configuração do ambiente e execução dos experimentos. \\
\texttt{notebooks/TC1/} & Diretório principal contendo os Jupyter Notebooks dos experimentos. \\
\hspace{1em}\texttt{MQO.ipynb} & Experimento de Regressão Linear por Mínimos Quadrados (MQ). \\
\hspace{1em}\texttt{perceptron-...ipynb} & Experimento com o modelo Perceptron de Regressão. \\
\hspace{1em}\texttt{MLP\_H1.ipynb} & Experimento com MLP de uma camada oculta. \\
\hspace{1em}\texttt{MLP\_H2.ipynb} & Experimento com MLP de duas camadas ocultas. \\
\texttt{trabalho\_ic\_aplicada/} & Pacote Python que contém a lógica dos modelos e funções de suporte. \\
\hspace{1em}\texttt{models/ridge\_reg....py} & Implementação da Regressão Linear (MQ). \\
\hspace{1em}\texttt{models/reg\_perc....py} & Implementação do Perceptron para Regressão. \\
\hspace{1em}\texttt{models/reg\_mlp.py} & Implementação da MLP para Regressão. \\
\hspace{1em}\texttt{models/aux.py} & Funções auxiliares para métricas (REQM, $R^2$), validação cruzada e geração de gráficos. \\
\texttt{results/TC1/} & Diretório destinado às saídas dos experimentos (histogramas de resíduos, gráficos de dispersão, etc.). \\
\bottomrule
\end{tabular}
\end{table}



\subsection{Ameaças à validade e controles}
Toda análise empírica está sujeita a limitações que podem comprometer a validade interna e externa dos resultados. No presente trabalho, destacamos algumas ameaças relevantes e as estratégias adotadas para mitigá-las, para assegurar maior robustez às conclusões.  

\begin{enumerate}
    \item \textbf{Tamanho amostral moderado} (\(n=414\)), sujeito à influência de pontos extremos.
    \item \textbf{Heterocedasticidade e caudas pesadas}, que podem afetar a confiabilidade das inferências.
    \item \textbf{Colinearidade entre atributos espaciais}, inflando potencialmente as incertezas em modelos lineares.
    \item \textbf{Particionamento aleatório}, que não controla a autocorrelação espacial existente nos dados.
\end{enumerate}

Como medidas de controle, evitamos vazamento nas normalizações, utilizamos validação cruzada no conjunto de treinamento, reportamos predições OOF agregadas para diagnóstico e mantivemos critérios uniformes de seleção de modelos.
